\chapter{Leggi congiunte di variabili aleatorie}
\section{Funzioni di distribuzione congiunte}

\begin{Definition}{Funzione di distribuzione congiunta}{}
    Date due variabili aleatorie (continue o discrete) $X$ e $Y$, definiamo la funzione di distribuzione congiunta la funzione di due variabili $a,b\in(-\infty,+\infty)$:
    \begin{equation*}
        F(a,b)=P(X\leq a, Y\leq b)
    \end{equation*}
\end{Definition}

\begin{Definition}{Funzione di densità discreta congiunta}{}
    Date due variabili aleatorie discrete $X$ e $Y$ con valori $I$ definiamo la funzione di densità discreta congiunta come:
    \begin{equation*}
        p(x,y)=P(X=x,Y=y) 
    \end{equation*}
    La densità discreta di $X$ può essere ottenuta da $p(x, y)$ tramite:
    \begin{equation*}
        p_X(x)=P(X=x)=\sum_{y:p(x,y)>0}p(x,y)
    \end{equation*}
    In maniera analoga:
    \begin{equation*}
        p_Y(y)=\sum_{x:p(x,y)>0}p(x,y)
    \end{equation*}
\end{Definition}

Nelle distribuzioni multinomiali abbiamo $n$ prove indipendenti identiche, dove ogni sigola prova ha $r\in\mathbb{N}$ possibili esiti, $E_1, E_2, ..., E_r$ con probabilità rispettivamente $p_1, p_2, ..., p_r$ e $\sum^r_{i=1}p_i=1$. Essendo $X_i$ il numero di prove il cui esito è $E_i$ la densità congiunta di $X_1, X_2, ..., X_r$ sarà:
\begin{equation*}
    P(X_1=n_1, X_2=n_2, ..., X_r=n_r)=\begin{cases}
        \frac{n!}{n_1!n_2!...n_r!}p_1^{n_1}p_2^{n_2}...p_r^{n_r}\hspace{1cm
        }\mbox{se }\sum^r_{i=i}n_i=n\\0\hspace{3.9cm}\mbox{se }\sum^r_{i=i}n_i\neq n
    \end{cases}
\end{equation*}

\begin{Example}{Distribuzioni multinomiali}{}
    \textit{Supponiamo di avere un'urna con 3 palline blu, 2 rosse e 3 verdi. Peschiamo 5 palline con reinserimento. Definiamo la funzione di densità congiunta di $X_1, X_2, X_3$.}

    Le prove sono identiche e si utilizza la distribuzione multinomiale con $n=5$ e $r=3$. Le probabilità sono $p_1=\frac{3}{8}, p_2=\frac{2}{8}, p_3=\frac{3}{8}$, abbiamo che:
    \begin{equation*}
        P(X_1=n_1,X_2=n_2,X_3=n_3)=\begin{cases}
            \binom{5}{n_1,n_2,n_3}\left(\frac{3}{8}\right)^{n_1}\left(\frac{2}{8}\right)^{n_2}\left(\frac{3}{8}\right)^{n_3}\hspace{1cm}\mbox{se }n_1+n_2+n_3=n\\0\hspace{4.9cm}\mbox{altrimenti}
        \end{cases}
    \end{equation*}
\end{Example}

\begin{Definition}{Densità congiunta di variabili aleatorie continue}{}
    Diciamo che le variabili aleatorie $X$ e $Y$ sono congiuntamente (assolutamente) continue se esiste una funzione $f(x, y)$ integrabile tale che:
    \begin{equation*}
        P((X,Y)\in C)=\int\int _{(x,y)\in C}f(x,y)dxdy
    \end{equation*}
    La funzione $f(x, y)$ è chiamata (funzione di) densità congiunta di $X$ e $Y$. Se $A$ e $B$ sono una qualsiasi coppia di sottoinsiemi della retta reale, allora definendo $C = \{(x, y):x\in A, y\in B\}$ otteniamo che:
    \begin{equation*}
        P(X\in A, Y\in B)=\int_B\int_Af(x,y)dxdy
    \end{equation*}
\end{Definition}

Se $X$ e $Y$ sono congiuntamente continue, esse lo saranno anche individualmente, e le loro densità si posso ottenere come segue:
\begin{equation*}
    P(X\in A)=P\{X\in A, Y\in(-\infty,\infty)\}=\int_A\int^\infty_{-\infty}f(x,y)dydx=\int_Af_X(x)dx
\end{equation*}
Possiamo così trovare la densità di $X$ e, in maniera analoga, la densità di $Y$:
\begin{equation*}
    f_X(x)=\int^\infty_{-\infty}f_X(x,y)dx\hspace{1cm}f_Y(y)=\int^\infty_{-\infty}f(x,y)dx
\end{equation*}


\begin{Example}{Densità congiunta di variabili aleatorie continue}{}
    \textit{La densità congiunta di $X$ e $Y$ è data da:}
    \begin{equation*}
        f(x,y)=\begin{cases}
            2e^{-x}e^{-2y}\hspace{1cm}0<x<\infty,0<y<\infty\\0\hspace{2.1cm}\mbox{altrimenti}
        \end{cases}
    \end{equation*}
    \textit{Determinare (a) $P(X>1,Y<1)$ e (b) $P(X<a)$.}

    \begin{enumerate}[label=(\alph*)]
        \item 
        \begin{align*}
            P(X>1,Y<1)=\int^1_0\int^\infty_12e^{-x}e^{2y}dxdy=\int^1_02e^{-2y}\left(-e^{-x}\Bigm\lvert^\infty_1\right)dy=e^{-1}\int^1_02e^{-2y}dy=\\e^{-1}(1-e^2)
        \end{align*}

        \item 
        \begin{align*}
            P(X<a)=\int^a_0\int^\infty_02e^{-2y}e^{-x}dydx=\int^a_0e{-x}dx=1-e^{-a}
        \end{align*}
    \end{enumerate}
\end{Example}

\section{Variabili aleatorie indipendenti}

\begin{Definition}{Variabili aleatorie indipendenti}{}
    Due variabili aleatorie $X$ e $Y$ (continue o discrete si dicono indipendenti se:
    \begin{equation*}
        P(X\in A, Y\in B)=P(X\in A)P(Y\in B)
    \end{equation*}
    In altre parole:
    \begin{equation*}
        P(X\in A, Y\in B)=P(X\leq a, Y\leq b)=P(X\leq a)P(Y\leq b)
    \end{equation*}
\end{Definition}

Quindi, in termini della distribuzione congiunta $F$ di $X$ e $Y$, abbiamo che $X$ e $Y$ sono indipendenti se e solo se:
\begin{equation*}
    F(a,b)=F_X(a)F_Y(b)
\end{equation*}

Quando X e Y sono variabili aleatorie discrete o continue, la definizione d’indipendenza risulta rispettivamente:
\begin{equation*}
    p(x,y)=p_X(x)p_Y(Y)\hspace{1cm}f_{X,Y}=f_X(x)f_Y(y)
\end{equation*}

\begin{Example}{Variabili aleatorie indipendenti}{}
    \textit{$X$, $Y$ sono due variabili aleatorie con valori in $I=\{0,1\}$ e densità congiunta $p_{x,y}(0,0)=p_{X,Y}(0,1)=p_{X,Y}(1,0)=p_{X,Y}(1,1)=\frac{1}{4}$. Stabilire se $X$ e $Y$ sono indipendenti.}

    Determiniamo la densità di $X$ e di $Y$:
    \begin{equation*}
        p_X(x)=\sum_{y\in I}p_{X,Y}(x,y)\hspace{1cm}p_Y(y)=\sum_{x\in I}p_{X,Y}(x,y)
    \end{equation*}
    \begin{align*}
        &p_X(0)=p_{X,Y}(0,0)+p_{X,Y}(0,1)=\frac{1}{4}+\frac{1}{4}=\frac{1}{2}\\
        &p_X(1)=p_{X,Y}(1,0)+p_{X,Y}(1,1)=\frac{1}{4}+\frac{1}{4}=\frac{1}{2}\\
        &p_Y(0)=p_{X,Y}(0,0)+p_{X,Y}(1,0)=\frac{1}{4}+\frac{1}{4}=\frac{1}{2}\\
        &p_Y(0)=p_{X,Y}(0,1)+p_{X,Y}(1,1)=\frac{1}{4}+\frac{1}{4}=\frac{1}{2}
    \end{align*}
    \begin{equation*}
        p_{X,Y}=\frac{1}{4}=p_X(x)p_Y(y)=\frac{1}{2}\frac{1}{2}=\frac{1}{4}
    \end{equation*}
    $X$ e $Y$ sono indipendenti.
\end{Example}

\begin{Definition}{Variabili aleatorie identicamente distribuite}{}
    Le variabili aleatorie $X_1, X_2, ...,X_n)$ si dicono identicamente distribuite ed indipendenti (si scrive i.i.d.) se valgono:
    \begin{enumerate}
        \item tutte le variabili aleatorie hanno la stessa densità (discreta o continua) e distribuzione;
        \item
        \begin{equation*}
            P(X_1\in A_1, X_2\in A_2, ...,X_n\in A_n)=P(X_1\in A_1)P(X_2\in A_2)...P(X_n\in A_n)
        \end{equation*}
    \end{enumerate}
\end{Definition}

\section{Covarianza, varianza e somma e correlazioni}
\begin{Definition}{Covarianza}{}
    La covarianza di $X$ e $Y$, indicata con $Cov(X, Y)$, è definita come:
    \begin{equation*}
        Cov(X,Y):=E((X-E(X))(Y-E(Y)))
    \end{equation*}
     o equivalentemente:
     \begin{equation*}
         Cov(X,Y)=E(XY)-E(X)E(Y)
     \end{equation*}
\end{Definition}

Le proprietà della covarianza sono:
\begin{enumerate}
    \item $Cov(X, Y) = Cov(Y, X)$
    \item $Cov(X, X) = Var(X)$
    \item $Cov(aX,Y) = a Cov(X, Y)$
    \item $Cov(\sum^n_{i=1}X_i,\sum^m_{j=1}Y_j)=\sum^n_{i=1}\sum^m_{j=1}Cov(X_i,Y_j)$
\end{enumerate}

Se $X$ e $Y$ sono indipendenti allora $Cov(X,Y)=0$, ma questo non implica che se $Cov(X,Y)=0$ allora $X$ e $Y$ sono indipendenti.